%!TEX root = ../graduate-work.tex

\section{Обзор существующих решений}

% Здесь надо рассмотреть все существующие решения поставленной задачи, но не просто пересказать, в чем там дело, а оценить степень их соответствия тем ограничениям, которые были сформулированы в постановке задачи.

% Информацию о группах автоморфизмов других известных кодов можно найти в~\cite{FirstArticle}.

На данный момент существует много способов оценки качества автоматически сгенерированного текста. В литературе по генерации текстов, включая радиологические отчёты, традиционно используются общие методы оценки из области генерации естественного языка (natural language generation, NLG), заимствованные из машинного перевода. К ним относятся BLEU, METEOR, ROUGE, CIDEr, SPICE и BERTScore~\cite{papineni2002bleu,banerjee2005meteor,lin2004rouge,vedantam2015cider,anderson2016spice,zhang2020bertscore}. Более поздние обзоры, посвященные медицинским текстам, дополняют их специализированными подходами на основе структурированных представлений содержимого отчёта: RadGraph F1, векторное представление CheXbert (SembScore), RadCliQ, RaTEScore, а также методы оценки на основе больших языковых моделей, такие как GREEN и FineRadScore~\cite{yu2023evaluating, smit-etal-2020-combining, jain2021radgraph,ostmeier2024green,pmlr-v252-huang24a, zhao2024ratescore}.

\subsection{Общие NLG-подходы}
\textbf{BLEU}. Bilingual Evaluation Understudy (BLEU) -- один из первых и наиболее распространенных до сих пор методов оценки качества машинного перевода~\cite{papineni2002bleu}. Он опирается на модифицированную точность N-грамм $p_n$ между отчётом-кандидатом и эталонными описаниями. Для каждого порядка N она определяется как
\[
p_n \;=\; \frac{\sum\limits_{g \in \mathcal{G}_n} \min\!\bigl(\mathrm{count}_{\mathrm{cand}}(g),\, \mathrm{count}_{\mathrm{ref}}(g)\bigr)}{\sum\limits_{g \in \mathcal{G}_n} \mathrm{count}_{\mathrm{cand}}(g)},
\]
где $\mathcal{G}_n$ -- множество всех N-грамм, встречающихся в кандидате; $\mathrm{count}_{\mathrm{cand}}(g)$ -- число вхождений N-граммы $g$ в кандидате, а $\mathrm{count}_{\mathrm{ref}}(g)$ -- <<обрезанная>> частота, равная числу вхождений $g$ в эталонный текст.

Итоговый балл вычисляется как взвешенное геометрическое среднее этих точностей с учётом штрафа за краткость:
\[
\mathrm{BLEU} = \mathrm{BP} \cdot \exp\left(\sum_{n=1}^{N} w_n \log p_n \right),
\]
где $\mathrm{BP}$ -- штраф за краткость, $\mathrm{BP}=1$ при $c>r$ и $\mathrm{BP}=\exp(1-r/c)$ иначе; $c$ и $r$ -- суммарные длины кандидатов и эталонов соответственно.

Данный подход работает очень быстро и применим к любому языку. Однако у него есть ряд ограничений: в методе BLEU учитываются только точные совпадения в формулировках, а не семантика текста, и поэтому два текста, одинаковых по смыслу, но сформулированных по-разному, могут получить низкую оценку. Для исправления этих недостатков были предложены более поздние методы.

\textbf{ROUGE}. Семейство методов ROUGE~\cite{lin2004rouge} было разработано для оценки систем автоматического реферирования. В изученных работах чаще всего применяется ROUGE-L, основанный на длине наибольшей общей подпоследовательности между эталонным и сгенерированным текстами. Пусть $LCS(X,Y)$ -- длина наибольшей общей подпоследовательности для двух последовательностей слов $X$ и $Y$, тогда
    \[
    P_{\text{LCS}} = \frac{LCS(X,Y)}{|X|}, \qquad
    R_{\text{LCS}} = \frac{LCS(X,Y)}{|Y|},
    \]
    где $|X|$ и $|Y|$ обозначают длину последовательностей (число слов в эталонном и сгенерированном тексте соответственно). Итоговый балл задается $F$-мерой:
    \[
    F_{\text{LCS}} = \frac{(1+\beta^2)P_{\text{LCS}}R_{\text{LCS}}}{R_{\text{LCS}} + \beta^2 P_{\text{LCS}}},
    \]
    где $\beta \geq 0$ -- гиперпараметр, задающий относительный вес полноты по сравнению с точностью. ROUGE учитывает полноту, а не только точность, в отличие от BLEU, однако также неспособен оценивать смысл текста.

\textbf{METEOR}. METEOR~\cite{banerjee2005meteor} также исходит из совпадений слов, но сначала строит выравнивание между сгенерированным и эталонным текстами (с учётом стемминга (нахождения корня слова) и синонимов), после чего вычисляет точность $P$ и полноту $R$ по униграммам. Гармоническое среднее задается с большим весом полноты:
    \[
    F_{\mathrm{mean}} = \frac{10PR}{R + 9P},
    \]
    а затем вводится штраф за фрагментацию $\mathrm{Penalty}$:
    \[
    \mathrm{Penalty} = 0.5 \left( \frac{\#\text{chunks}}{\#\text{unigrams\_matched}} \right)^3,
    \]
    зависящий от числа непрерывных последовательностей совпавших слов (chunks). 
    
    Итоговое значение определяется как
    \[
    \mathrm{METEOR} = F_{\mathrm{mean}} \cdot (1 - \mathrm{Penalty}).
    \]
    
    Данный метод лучше коррелирует с человеческой оценкой, чем рассмотренные ранее методы, однако все ещё неспособен в полной мере учесть смысловое наполнение текстов.


\textbf{CIDEr}.
Подход CIDEr (Consensus-based Image Description Evaluation) был предложен для оценки описаний изображений (image captioning) и измеряет степень согласованности описания-кандидата с набором эталонных описаний одного и того же изображения~\cite{vedantam2015cider}. Для повышения чувствительности к содержательно важным словам используется TF–IDF-взвешивание N-грамм; наиболее частые, малоинформативные N-граммы получают меньший вес. 

Рассмотрим изображение $i$ с описанием-кандидатом $c_i$ и набором эталонных описаний $S_i = \{s_{ij}\}_{j=1}^{m}$. После стемминга для каждого $N \in \{1,2,3,4\}$ каждое предложение представляется набором $N$-грамм $\{w_k\}$. Пусть $h_k(s_{ij})$ и $h_k(c_i)$ обозначают количество появлений N-граммы $w_k$ в $s_{ij}$ и $c_i$ соответственно. Тогда для предложения $s_{ij}$ определим компоненту TF–IDF для N-граммы $w_k$ как
\[
g_k(s_{ij}) = \underbrace{\frac{h_k(s_{ij})}{\sum\limits_{w_l} h_l(s_{ij})}}_{\text{TF}} \cdot
\underbrace{\log \frac{|I|}{\sum\limits_{p=1}^{|I|} \mathbb{I}\Bigl[\sum\limits_{j} h_k(s_{pj}) > 0\Bigr]}}_{\text{IDF}},
\]
где $|I|$ -- количество изображений в корпусе, а индикатор $\mathbb{I}[\cdot]$ показывает, встречается ли N-грамма $w_k$ хотя бы в одном эталонном описании изображения $p$. Аналогично вычисляется $g_k(c_i)$ для описания-кандидата. Затем для данного $n$ вычисляется среднее косинусное сходство кандидата с эталонными описаниями:

%Для фиксированной длины N-грамм $n$ формируется TF–IDF-вектор
%\[
%\mathbf{g}^{(n)}(s_{ij}) = \bigl(g_k(s_{ij})\bigr)*{w_k \in \Omega_n},
%\qquad
%\mathbf{g}^{(n)}(c_i) = \bigl(g_k(c_i)\bigr)*{w_k \in \Omega_n},
%\]
%где $\Omega_n$ -- словарь всех N-грамм длины $n$. Затем для данного $n$ вычисляется среднее косинусное сходство кандидата с эталонными описаниями:
\[
\mathrm{CIDEr}_n(c_i, S_i) = \frac{1}{m} \sum_{j=1}^m
\frac{ g_n(c_i) \cdot g_n(s_{ij}) }{ \lVert g_n(c_i)\rVert \lVert g_n(s_{ij})\rVert }.
\]
Итоговый балл CIDEr для изображения $i$ определяется как взвешенная сумма по различным длинам N-грамм:
\[
\mathrm{CIDEr}(c_i, S_i) = \sum_{n=1}^{N} w_n \mathrm{CIDEr}_n(c_i, S_i),
\]
где в исходной работе обычно берут $N = 4$ и равные веса $w_n$.

Использование TF–IDF позволяет понижать вклад часто встречающихся N-грамм и усиливать вклад более информативных выражений. Однако для хорошей работы данный подход требует большого количества эталонных текстов, а также не полностью учитывает семантическую близость текстов.

\textbf{SPICE}. Метод оценки SPICE~\cite{anderson2016spice} также применяется в области image captioning. Он оценивает семантическое содержимое текстов: кандидатное и эталонные описания переводятся в графы сцены, в которых узлы соответствуют объектам, атрибутам и отношениям между ними. Из этих графов извлекаются множества логических кортежей $T(c)$ и $T(r)$, далее рассчитываются точность и полнота:
\[
P = \frac{|T(c)\cap T(r)|}{|T(c)|}, \quad
R = \frac{|T(c)\cap T(r)|}{|T(r)|},
\]
а итоговый балл -- стандартная $F$-мера:
\[
\mathrm{SPICE}(c,r) = \frac{2PR}{P+R}.
\]

Недостатком этого метода является то, что для корректного его использования необходима грамматическая и синтаксическая корректность оцениваемых предложений.

\textbf{BERTScore}. Подход BERTScore~\cite{zhang2020bertscore} использует векторные представления токенов, полученных из предобученной языковой модели BERT. Пусть $x=\langle x_1,\dots,x_m\rangle$ и $y=\langle y_1,\dots,y_n\rangle$ -- эталон и кандидат, а $\mathbf{r}_i$ и $\mathbf{c}_j$ -- их векторные представления. Определим матрицу косинусных сходств $s_{ij} = \cos(\mathbf{c}_j,\mathbf{r}_i)$. Тогда
\[
P = \frac{1}{n} \sum_{j=1}^n \max_i s_{ij}, \qquad
R = \frac{1}{m} \sum_{i=1}^m \max_j s_{ij},
\]
и итоговая оценка -- $F$-мера
\[
F_{\text{BERT}} = \frac{2PR}{P+R}.
\]
При необходимости значения могут дополнительно взвешиваться IDF-весами токенов для снижения влияния частых слов.

Все перечисленные NLG-подходы хорошо отражают лексическое и, частично, семантическое сходство текстов, но слабо чувствительны к клинически значимым расхождениям (например, наличию отрицания), что ограничивает их полезность для оценки радиологических отчётов.

\subsection{Подходы к оценке медицинских отчётов}

Далее будут рассмотрены методы оценки, специально разработанные для радиологических отчётов; именно они лежат в основе современных публичных лидербордов (рейтинговых таблиц) по генерации отчётов, таких как ReXrank~\cite{zhang2024rexrank}, и используются для объективного сравнения существующих моделей.

\textbf{RadGraph F1}. RadGraph~\cite{jain2021radgraph} предлагает схему разметки отчётов в виде графа клинических сущностей (патологии, анатомические структуры, модификаторы) и отношений между ними. Оценка RadGraph F1 вычисляется как $F$-мера между множествами сущностей и отношений, извлечённых из отчёта-кандидата и эталонного отчёта:
\[
F_{\text{RadGraph}} = \frac{2P_{\text{ent/rel}} R_{\text{ent/rel}}}{P_{\text{ent/rel}} + R_{\text{ent/rel}}},
\]
где $P_{\text{ent/rel}}$ и $R_{\text{ent/rel}}$ -- точность и полнота по сущностям и отношениям. Такой подход напрямую сравнивает факт наличия и локализацию патологий, но опирается на точность алгоритма извлечения графа.

\textbf{SembScore (косинусное сходство векторов CheXbert)}. В ра\-боте CheX\-pert~\cite{smit-etal-2020-combining} отчёту ставится в соответствие бинарный вектор меток патологий $y\in\{0,1\}^K$, кодирующий наличие либо отсутствие фиксированного набора находок. Сходство между отчётом-кандидатом и эталонным отчётом оценивается по косинусному сходству их векторных представлений:
\[
\mathrm{SembScore}\bigl(y^{(\text{ref})}, y^{(\text{cand})}\bigr) =
\frac{ y^{(\text{ref})} \cdot y^{(\text{cand})} }{ \lVert y^{(\text{ref})}\rVert \lVert y^{(\text{cand})}\rVert }.
\]
Этот подход хорошо отражает согласие по фиксированному набору диагнозов, но плохо покрывает редкие патологии и детали локализации. 

\textbf{RaTEScore}. RaTEScore~\cite{zhao2024ratescore} -- трёхэтапный подход: сначала разработанная авторами модель RaTE-NER выделяет из отчёта клинические сущности $e_i = (n_i, t_i)$ с типами из множества Anatomy (анатомия), Abnormality (аномалия), Disease (заболевание), Non-Abnormality (отсутствие аномалии), Non-Disease (отсутствие заболевания). Затем названия сущностей $n_i$ отображаются кодировщиком BioLORD~\cite{remy2024biolord} в векторные представления $f_i$, так что синонимы оказываются близкими в пространстве векторных представлений. На финальном этапе для каждой сущности кандидата выбирается наиболее близкая (по косинусному сходству) сущность эталона, совокупное сходство $S(x,\hat{x})$ вычисляется с учётом матрицы весов $W$ по типам сущностей и штрафа $p$ за несовпадение типов, после чего из величин $S(x,\hat{x})$ и $S(\hat{x},x)$, интерпретируемых как аналоги точности и полноты, получается итоговый RaTEScore в форме F-меры. Модель NER обучается на наборе данных RaTE-NER, кодировщик BioLORD используется предобученный, а параметры $W$ и $p$ находятся по размеченному врачами набору RaTE-Eval, представленному в статье авторов метода. 

Однако недостатком данного подхода является то, что он не учитывает структурные связи между сущностями, обрабатывая их как независимые элементы. Ещё одно существенное требование -- необходимость совпадения числа предложений в эталонном и сгенерированном отчётах. Кроме того, RaTEScore часто недооценивает степень сходства отчётов и выдаёт заниженные значения итогового счёта
~\cite{gupta2024analyzing}.

\textbf{RadCliQ}. RadCliQ~\cite{yu2023evaluating} -- обучаемый регрессионный подход, который совмещает 4 ранее описанных метода (BLEU, BERTScore, SembScore, RadGraph F1) в единый прогноз, приближающий оценку качества отчёта по числу ошибок, выявленных радиологами. В публичной рейтинговой таблице моделей генерации радиологических отчётов ReXrank~\cite{zhang2024rexrank} данный подход используется в качестве основой оценки каждой из сравниваемых моделей.

\textbf{CREPE.}
Метод CREPE~\cite{cho2025crepe} предсказывает число ошибок в нескольких категориях. На вход подаются эталонный и кандидатный отчёты, после чего пара текстов переводится в векторное представление с помощью специализированного кодировщика BiomedBERT. Поверх общего текстового представления используются несколько регрессионных голов, каждая из которых предсказывает число ошибок определённого типа. В работе рассматриваются шесть категорий ошибок: ложное указание находки, пропуск находки, неверная локализация, неверная степень выраженности, ошибочное упоминание сравнения и пропуск сравнения. Итоговая оценка CREPE получается суммированием предсказанных значений по категориям, поэтому меньшая оценка соответствует более качественному заключению.

\subsection{Подходы к оценке на основе больших языковых моделей}

\textbf{GREEN}. GREEN (Generative Radiology Report Evaluation and Error Notation)~\cite{ostmeier2024green} использует большую языковую модель для выделения и подсчёта клинически значимых и незначимых ошибок в отчёте-кандидате относительно эталонного отчёта. Модель генерирует структурированное описание ошибок (ложные находки, пропущенные находки, ошибки локализации, ошибки степени выраженности и др.) и агрегированный балл каждой пары отчётов. В упрощенном виде эту оценку можно интерпретировать как нормализованную долю корректных клинических утверждений:
\[
\mathrm{GREEN} = \frac{N_{\text{matched findings}}}{N_{\text{matched findings}} + \sum_i N_{\text{error}_i}},
\]
где $N_{\text{error}_i}$ -- число клинически значимых ошибок различных типов, выделенных моделью. Типы ошибок, выделяемых моделью:
\begin{enumerate}
    \item Ложное указание находки (находка упомянута, но в эталоне её нет).
    \item Пропуск находки (находка есть в эталоне, но отсутствует в кандидате).
    \item Неверная локализация/расположение находки.
    \item Неверная степень выраженности (тяжесть) находки.
    \item Упоминание сравнения, которого нет в эталонном заключении.
    \item Пропуск сравнения, описывающего изменения относительно предыдущего исследования.
\end{enumerate}

Модель GREEN была получена путём дистилляции модели GPT-4 и содержит менее 7 миллиардов параметров, что позволяет запускать ее локально. Метод демонстрирует более высокую корреляцию с оценками радиологов по сравнению с BLEU, CIDEr, RadGraph F1 и RadCliQ, а также содержит интерпретируемые текстовые объяснения выделенных ошибок.

\textbf{FineRadScore}. FineRadScore~\cite{pmlr-v252-huang24a} формулирует задачу оценки как поиск минимального числа строковых правок, необходимых для преобразования отчёта-кандидата в эталонный отчёт, с использованием большой языковой модели. Для каждой строки модель генерирует исправленный вариант, оценку серьёзности ошибки с медицинской точки зрения и поясняющий комментарий. Число правок и их серьёзность затем агрегируются в итоговый балл, убывающий при увеличении количества и тяжести выявленных ошибок. Авторы показывают, что FineRadScore хорошо согласуется с оценками радиологов как на уровне отдельных строк, так и на уровне отчёта в целом.

\textbf{GEMA-Score.} GEMA-Score (Granular Explainable Multi-Agent Score)~\cite{zhang2025gema}
представляет собой многоагентную систему на основе больших языковых моделей. Метод сравнивает отчёт-кандидат и отчёт-эталон по клинически значимым аспектам, включая наличие находок, локализацию, степень выраженности и неопределённость, а также учитывает более общие характеристики текста, такие как полнота, читаемость и клиническая применимость. Итоговая оценка формируется на основе частных решений агентов и сопровождается интерпретируемым объяснением выявленных расхождений.

Согласно результатам работ~\cite{yu2023evaluating,ostmeier2024green,pmlr-v252-huang24a,zhao2024ratescore}, наилучшее согласие с экспертными оценками в задаче оценки радиологических отчётов обычно демонстрируют подходы на основе больших языковых моделей, в частности GREEN и FineRadScore.
